\documentclass[11pt]{article}

% cs250preamble.tex --- shared preamble for CS 250 course notes
%
% This file is \input by main.tex and by each individual module
% file so that each module can still compile standalone. Any
% package or custom-environment change should be made here, not
% in the individual module files.

\usepackage[T1]{fontenc}
\usepackage[margin=1in]{geometry}
\usepackage{amsmath, amssymb, amsthm}
\usepackage{enumitem}
\usepackage{hyperref}
\usepackage{booktabs}
\usepackage{listings}
\usepackage{xcolor}
\usepackage{tikz}
\usetikzlibrary{shapes.geometric, shapes.gates.logic.US, shapes.gates.logic.IEC, arrows.meta, positioning, calc, trees}
\usepackage{bussproofs}
\usepackage{mdframed}

\lstset{
  basicstyle=\ttfamily\small,
  frame=single,
  breaklines=true,
  columns=fullflexible,
  mathescape=true
}

\theoremstyle{definition}
\newtheorem{theorem}{Theorem}[section]
\newtheorem{definition}[theorem]{Definition}
\newtheorem{example}[theorem]{Example}

\hypersetup{
  colorlinks=true,
  linkcolor=blue!60!black,
  urlcolor=blue!60!black
}

% Key-takeaway box: used at the end of major sections and at the end of
% each module to summarise the main points. Expects \item bullets inside.
\newmdenv[
  linewidth=0.8pt,
  linecolor=blue!40!black,
  backgroundcolor=blue!3,
  skipabove=8pt,
  skipbelow=8pt,
  innertopmargin=7pt,
  innerbottommargin=7pt,
  innerleftmargin=9pt,
  innerrightmargin=9pt
]{keytakeawaybox}
\newenvironment{keytakeaway}{%
  \begin{keytakeawaybox}%
  \noindent\textbf{Key takeaways.}%
  \begin{itemize}[leftmargin=1.3em, topsep=2pt, itemsep=1pt, label=\textbullet]%
}{%
  \end{itemize}%
  \end{keytakeawaybox}%
}

% Summary/looking-ahead environment: used once at the end of each module
% to wrap up what the module built and point forward.
\newmdenv[
  linewidth=0.8pt,
  linecolor=gray!60,
  backgroundcolor=gray!5,
  skipabove=8pt,
  skipbelow=8pt,
  innertopmargin=7pt,
  innerbottommargin=7pt,
  innerleftmargin=9pt,
  innerrightmargin=9pt
]{summarybox}

% Sidebar environment: used for tangential asides---historical notes,
% pointers to other areas of math/CS, cross-paradigm remarks---that
% enrich the main narrative without being essential to it. Takes one
% argument: the sidebar's title.
\newmdenv[
  linewidth=0.6pt,
  linecolor=gray!60,
  backgroundcolor=gray!4,
  skipabove=8pt,
  skipbelow=8pt,
  innertopmargin=6pt,
  innerbottommargin=6pt,
  innerleftmargin=9pt,
  innerrightmargin=9pt
]{sidebarbox}
\newenvironment{sidebar}[1]{%
  \begin{sidebarbox}%
  \noindent\textbf{Sidebar: #1.}\par\medskip%
}{%
  \end{sidebarbox}%
}

% Reading-map environment: used at the top of each numbered module to
% indicate Epp coverage, prerequisites, relevant intermezzi, and
% suggested practice problems. Expects four \rmentry{label}{body}
% items inside (or arbitrary paragraphs).
\newmdenv[
  linewidth=0.8pt,
  linecolor=black!55,
  backgroundcolor=yellow!4,
  skipabove=10pt,
  skipbelow=14pt,
  innertopmargin=9pt,
  innerbottommargin=9pt,
  innerleftmargin=11pt,
  innerrightmargin=11pt
]{readingmapbox}
\newcommand{\rmentry}[2]{\par\noindent\textbf{#1}\hspace{0.4em}#2\par}
\newenvironment{readingmap}{%
  \begin{readingmapbox}%
  \noindent{\bfseries\large Reading map}%
  \vspace{2pt}%
  \setlength{\parskip}{4pt plus 1pt}%
}{%
  \end{readingmapbox}%
}


\title{CS 250 --- Module 2: Propositional Logic}
\author{}
\date{}

\begin{document}
\maketitle

\begin{readingmap}
\rmentry{Epp sections.}{§2.1--2.2: logical form and logical equivalence.}

\rmentry{Prerequisites.}{Module 1 (sets, functions, truth as a boolean-valued assignment).}

\rmentry{Relevant intermezzi.}{The Gentzen/inference-tree intermezzo (optional, useful preview for Module 3).}

\rmentry{Practice from Epp.}{§2.1 exercises on building truth tables and classifying formulas as tautology / contradiction / contingent; §2.2 exercises on logical equivalences and De Morgan's laws.}
\end{readingmap}

\section{Logical formulas, formalized}
With all our mathematical preliminaries out of the way, we're ready to start dealing with logic more formally. Early attempts at talking about logic were focused on rhetoric and how to make arguments in human language. They were still about reasoning but didn't treat it in the \emph{symbolic} way we do in modern times, where we turn logic into formulas and equations and symbols.

In this class, I'm not only going to take a symbolic approach to logic like Epp but I'm going to take a very ``computer sciencey'' approach, discussing logic as though it were a programming language.

There are many possible logics --- yes, plural, because there are many possible logics just as there are many possible programming languages, each serving its own purpose. The standard way to introduce any logic is to define its \emph{syntax} and its \emph{semantics}.

The \emph{syntax}\index{syntax} of the logic is a lot like when we talk about the syntax of any programming language: what combinations of symbols make a ``valid'' program/formula. The syntax tells us how we're allowed to build formulas in our logic.

The \emph{semantics}\index{semantics} of a logic tells us what the formulas of the logic \textbf{mean}. In ordinary programming languages, the compiler or interpreter implements the semantics: it defines what a syntactically valid program does when executed. For (most) logics, the semantics tells us how to take a formula and turn it into a mathematical function that takes in arguments and returns either ``true'' or ``false''.

What do true and false mean here? That question could fill a PhD thesis, but for our purposes: we can define ``true'' as ``reflects reality'' and ``false'' as ``contradicts reality''. For most of this class, you can think of ``true'' and ``false'' a lot like the boolean types you find in every programming language.

The first logic we're going to deal with is really simple. It only has booleans as data, variables, and a few logical connectives: and ($\wedge$), or ($\vee$), not ($\sim$), and implication ($\to$). (We'll use $\lnot$ for negation in our mathematical notation; some textbooks, including Epp, write $\sim$ instead. They mean the same thing.) The idea of this logic is that it encompasses basic operations for how to connect statements about the world and their relationship to each other. It's all verbs, no nouns, but it still helps give some intuition of how reasoning works.

We're going to express this using something called Backus-Naur form\index{BNF (Backus-Naur form)} (\url{https://en.wikipedia.org/wiki/Backus-Naur_form}) which is a standard notation for describing grammars of formal languages, like this one, and something you'll see more of as you study computer science.

\begin{definition}[Propositional Logic Syntax]
The syntax of propositional logic\index{propositional logic} is given by the following BNF grammar:
\begin{lstlisting}
L := T | F | L ^ L | L v L | ~L | L -> L | x
\end{lstlisting}
\end{definition}

BNF is a way to describe how to build things recursively. It tells us, in this case, that to build an ``L'' you have a bunch of possibilities separated by the | symbol. You can follow this recursively, continuing to substitute in things for any L symbols that appear until you only have base cases (here: variables, T, and F).

When a formula is ambiguous---like \texttt{p \^{} q v r}, which could be read as $(p \wedge q) \vee r$ or $p \wedge (q \vee r)$---we follow the standard precedence: $\lnot$ binds tightest, then $\wedge$, then $\vee$, then $\to$. We use parentheses to override precedence when needed. This is similar to how \texttt{*} binds tighter than \texttt{+} in arithmetic.

Here are some examples to help clarify how this is used. This grammar tells us that

\begin{itemize}
  \item T is a valid formula
  \item F is a valid formula
  \item $T \wedge (x \vee y)$ is a valid formula
  \item $T \to L$ is not a valid formula because there's still an L
  \item T + F is not a valid formula because + isn't an allowed symbol in the syntax
  \item heck is a valid formula, a single variable named heck
\end{itemize}

That's the syntax. We need to define the \emph{semantics} of this logic. We've already named the pieces, so let's give an informal semantics before presenting the formal one:

\begin{itemize}
  \item T is just T, the boolean value for true
  \item F is just F, the boolean value for false
  \item $p \wedge q$ is \emph{and}\index{conjunction@$\wedge$ (conjunction)} and is true when p and q are both true
  \item $p \vee q$ is \emph{or}\index{disjunction@$\vee$ (disjunction)} and is true when \textbf{either} p or q are true, but \textbf{both} can be true
  \item $\lnot p$ is \emph{not}\index{negation@$\lnot$ (negation)} and is true when p isn't
  \item $p \to q$ is \emph{implication}\index{implication@$\to$ (implication)} and captures the idea that p being true forces q to be true. Specifically, it's false only when p is true and q is false
  \item variables are true when you plug in a true value into them
\end{itemize}

That's the informal semantics. What's the \emph{formal} semantics? We're going to explain it in terms of \textbf{truth tables}\index{truth table}.

A truth table is a way of describing the \emph{graph} of the function that a formula represents. Since all the arguments to these functions are just booleans and can be enumerated easily, the graph can be represented as a table.

Let's start with variables.

\begin{center}
\begin{tabular}{cc}
\toprule
$x$ & $x$ \\
\midrule
T & T \\
F & F \\
\bottomrule
\end{tabular}
\end{center}

A variable is just an identity function---it returns its input unchanged.

From here, we can build up the meaning of any formula piecemeal and describe the meaning of connectives in terms of ``p'' and ``q'' as subformulas, so when we write something like this:

\begin{center}
\begin{tabular}{cc|c}
\toprule
$p$ & $q$ & $p \wedge q$ \\
\midrule
T & T & T \\
T & F & F \\
F & T & F \\
F & F & F \\
\bottomrule
\end{tabular}
\end{center}

what we're saying is that if we've already evaluated the meaning of ``p'' and ``q'' we can plug in their values and get the meaning of when they're connected by an \emph{and}. So we can read this table off and it matches the informal meaning we described above: $p \wedge q$ is true exactly when p is true and q is true and false otherwise.

Next, $p \vee q$:

\begin{center}
\begin{tabular}{cc|c}
\toprule
$p$ & $q$ & $p \vee q$ \\
\midrule
T & T & T \\
T & F & T \\
F & T & T \\
F & F & F \\
\bottomrule
\end{tabular}
\end{center}

Which tells us that $p \vee q$ is true when either p or q is true (including both of them).

Next, negation:

\begin{center}
\begin{tabular}{c|c}
\toprule
$p$ & $\lnot p$ \\
\midrule
T & F \\
F & T \\
\bottomrule
\end{tabular}
\end{center}

Now we get to the counterintuitive one: implication. The intuition is that we want $p \to q$ to be true if p being true makes q be true---in other words, I can take the knowledge that p is true and then conclude that q is true.

(Sidenote: there's a sense, which we might get into later in this course, in which implication really is a function that takes in a proof that p is true and turns it \emph{into} a proof that q is true. Keep that idea in mind.)

\emph{Do I care} if q is true without p being true? In our first choice of logic, the answer is ``no''. We're going to make the decision that $F \to p$ is true no matter what p is, because the case where the premise is false is irrelevant to what we want implication to capture. This has some odd implications! From a false premise you can conclude anything. We'll dig into that more as we go.

For now, the truth table:

\begin{center}
\begin{tabular}{cc|c}
\toprule
$p$ & $q$ & $p \to q$ \\
\midrule
T & T & T \\
T & F & F \\
F & T & T \\
F & F & T \\
\bottomrule
\end{tabular}
\end{center}

With these truth tables in hand, let's show how to combine them together in an example:

\begin{center}
\begin{tabular}{ccc|l}
\toprule
$x$ & $y$ & $z$ & $((x \vee y) \wedge z) \to x$ \\
\midrule
T & T & T & $((T \vee T) \wedge T) \to T \Longrightarrow T$ \\
T & T & F & $((T \vee T) \wedge F) \to T \Longrightarrow F \to T \Longrightarrow T$ \\
T & F & T & $((T \vee F) \wedge T) \to T \Longrightarrow T$ \\
T & F & F & $((T \vee F) \wedge F) \to T \Longrightarrow F \to T \Longrightarrow T$ \\
F & T & T & $((F \vee T) \wedge T) \to F \Longrightarrow T \to F \Longrightarrow F$ \\
F & T & F & $((F \vee T) \wedge F) \to F \Longrightarrow F \to F \Longrightarrow T$ \\
F & F & T & $((F \vee F) \wedge T) \to F \Longrightarrow F \to F \Longrightarrow T$ \\
F & F & F & $F \to F \Longrightarrow T$ \\
\bottomrule
\end{tabular}
\end{center}

Whew! That was a lot, but hopefully you can see that, in principle, you could use truth tables to calculate the meaning of any formula as a function.

\section{Tautologies and equivalences}
Having talked about our very first logic, which is usually called ``propositional logic'', we can start discussing some of its implications and structure.

\begin{definition}[Tautology]
A \emph{tautology}\index{tautology} is a formula whose meaning is \textbf{always} true, no matter what values you put into all the variables.
\end{definition}

For example, for any formula p we have that $p \vee \lnot p$ is a tautology! Let's check the truth table:

\begin{center}
\begin{tabular}{c|l}
\toprule
$p$ & $p \vee \lnot p$ \\
\midrule
T & $T \vee F = T$ \\
F & $F \vee T = T$ \\
\bottomrule
\end{tabular}
\end{center}

No matter what value the sub-formula p takes, you get true.

\begin{definition}[Contradiction]
A \emph{contradiction}\index{contradiction} is a formula whose meaning is \textbf{always} false, no matter what values you put into the variables.
\end{definition}

For example, $p \wedge \lnot p$ is a contradiction---there's no assignment that makes it true. Contradictions are the dual of tautologies and will play a key role when we discuss proof by contradiction in the next module.

Since propositional logic is a really simple logic, one that only has booleans as data, the tautologies of propositional logic aren't individually profound, but they tell us important things about valid reasoning. The above tautology tells us that for any formula---for any statement about the world---either the claim is true or it is false.

Maybe that doesn't seem profound, but this is called ``the law of the excluded middle'' and is actually pretty controversial in the history of logic! In \emph{this} propositional logic and \emph{this} notion of ``truth'', however, we can always conclude that if something isn't false it must be true.

Why controversial? There's a tradition in logic---called \emph{constructive} or \emph{intuitionistic} logic---that refuses to accept $p \vee \lnot p$ as an axiom. The idea is that to prove ``$A$ or $B$'' you should have to actually show \emph{which one} is true, not just argue that one of them must be.

From a programming perspective this makes a lot of sense: a constructive proof of ``$A$ or $B$'' is like a program that actually \emph{computes} which case holds and hands you the evidence. Classical logic, the kind we're using, lets you say ``well, one of them \emph{has} to be true'' without ever telling you which---you get an existence result but not an algorithm. The connection between proofs and programs turns out to be deep and beautiful, and it's one of the reasons logic matters so much to computer science.

In this course we'll stick with classical logic (with excluded middle) because that's what Epp uses and what most intro courses assume. Fair warning: your instructor has a constructive bias, and it will probably peek through from time to time.

We'll come back to tautologies a bit later when we want to talk about reasoning and proof a little more explicitly, but for now just keep in mind the idea that they provide templates for what valid reasoning looks like.

Another thing we can do with truth tables is calculate \emph{equivalences}---we can figure out when formulas are equivalent to each other.

\begin{definition}[Logical Equivalence]
Two formulas are \emph{logically equivalent}\index{logical equivalence} exactly when they have the same truth table.
\end{definition}

Since we're defining the meaning of a formula \textbf{as its truth table}, this is a natural notion of equivalence. Put differently: logical equivalence is the right notion of sameness for propositional formulas. Two formulas with the same truth table may look wildly different on paper---$(p \to q)$ and $(\lnot p \vee q)$ are the canonical example---but nothing we can say \emph{about} them inside propositional logic can tell them apart. This is a recurring pattern: each kind of mathematical object comes with its own notion of ``the same,'' chosen so that it abstracts away the differences that don't matter. The intermezzo after Module~6 is the essay on that pattern.

We can also express equivalence \emph{inside} the logic by defining a shorthand: $p \leftrightarrow q$ means $(p \to q) \wedge (q \to p)$. This is called the \emph{biconditional}\index{biconditional@$\leftrightarrow$ (biconditional)} and reads ``if and only if.'' It is \textbf{not} a primitive connective in our grammar---it's just an abbreviation for a conjunction of two implications. You can compute its truth table by expanding the definition:

\begin{center}
\begin{tabular}{cc|c}
\toprule
$p$ & $q$ & $p \leftrightarrow q \;\;(\text{i.e.}\; (p \to q) \wedge (q \to p))$ \\
\midrule
T & T & T \\
T & F & F \\
F & T & F \\
F & F & T \\
\bottomrule
\end{tabular}
\end{center}

So the biconditional is true exactly when $p$ and $q$ have the same truth value---you can think of it like an equality test on booleans. Saying ``if and only if'' is the same as saying the implication goes both ways.

As an example, we're going to prove two formulas that end up being useful in some optimizations: de Morgan's laws\index{De Morgan's laws}! (\url{https://en.wikipedia.org/wiki/De_Morgan%27s_laws})

Let's start with the law that $\lnot(p \wedge q) \equiv (\lnot p) \vee (\lnot q)$. We'll calculate both of their truth tables and compare!

\begin{center}
\begin{tabular}{cc|c}
\toprule
$p$ & $q$ & $\lnot(p \wedge q)$ \\
\midrule
T & T & F \\
T & F & T \\
F & T & T \\
F & F & T \\
\bottomrule
\end{tabular}
\end{center}

and

\begin{center}
\begin{tabular}{cc|c}
\toprule
$p$ & $q$ & $(\lnot p) \vee (\lnot q)$ \\
\midrule
T & T & F \\
T & F & T \\
F & T & T \\
F & F & T \\
\bottomrule
\end{tabular}
\end{center}

These are clearly the same!

De Morgan's laws tell us how negation interacts with $\wedge$ and $\vee$: pushing $\lnot$ inside a conjunction or disjunction swaps one connective for the other.

This raises a natural question: how few operations do you actually need? This is where we introduce the concept of ``functional completeness''\index{functional completeness}: how few operations do you need to recreate all the logical operations on booleans?

For example, we have the operation $p \to q$. To recall, this operation has the following truth table:

\begin{center}
\begin{tabular}{cc|c}
\toprule
$p$ & $q$ & $p \to q$ \\
\midrule
T & T & T \\
T & F & F \\
F & T & T \\
F & F & T \\
\bottomrule
\end{tabular}
\end{center}

Now it turns out this is identical to the truth table of $(\lnot p) \vee q$

\begin{center}
\begin{tabular}{cc|l}
\toprule
$p$ & $q$ & $(\lnot p) \vee q$ \\
\midrule
T & T & $F \vee T = T$ \\
T & F & $F \vee F = F$ \\
F & T & $T \vee T = T$ \\
F & F & $T \vee F = T$ \\
\bottomrule
\end{tabular}
\end{center}

The truth tables match! Technically, we don't even need implication in this particular logic---though this equivalence is specific to classical propositional logic and does not hold in general if we have $\lnot$ and $\vee$.

This means that just having $\vee$, $\wedge$, and $\lnot$ is ``functionally complete'' because it can describe all the operations of this logic.

\section{Addendum: nand, xor, and how many binary operations are there?}
There are some additional operations that you \emph{can} define that we haven't defined yet.

The first of these we'll call xor\index{XOR} (sometimes shown with the $\oplus$ symbol), the \emph{exclusive} or and it has the following truth table:

\begin{center}
\begin{tabular}{cc|c}
\toprule
$p$ & $q$ & $p \oplus q$ \\
\midrule
T & T & F \\
T & F & T \\
F & T & T \\
F & F & F \\
\bottomrule
\end{tabular}
\end{center}

This is true if \emph{either} p or q is true but not when \emph{both} are true.

The second of these is \emph{nand}\index{NAND}, which is just ``not-and'' and it's rather important to the design of circuits:

\begin{center}
\begin{tabular}{cc|c}
\toprule
$p$ & $q$ & $p \mathbin{\text{nand}} q$ \\
\midrule
T & T & F \\
T & F & T \\
F & T & T \\
F & F & T \\
\bottomrule
\end{tabular}
\end{center}

An obvious question, and one that comes up in the homework of the accompanying class to these notes, is \emph{how many of these operations are there and what are they?}

We can approach this \emph{exhaustively}. A binary operation takes two (2) arguments, each of which is a boolean and can take two values. That's where we get the form of the truth tables above!

\begin{center}
\begin{tabular}{cc|c}
\toprule
$p$ & $q$ & $f(p,q)$ \\
\midrule
T & T & ? \\
T & F & ? \\
F & T & ? \\
F & F & ? \\
\bottomrule
\end{tabular}
\end{center}

How many different binary operations are even definable? That's equivalent to asking how many different truth tables are possible. There's a distinct truth table for every choice of how we fill in those four output entries in the table above. Well, okay, we know how to do this! Each of those four slots involves making a choice between one of two things: T and F. So the total number of choices must be $2 \times 2 \times 2 \times 2 = 2^4$ which is 16.

To prove \emph{functional completeness} of a set of boolean operators, you need to show that for any possible truth table you can find a way to replicate it using those operators. Let's do a concrete example:

Here's our truth table for implication

\begin{center}
\begin{tabular}{cc|c}
\toprule
$p$ & $q$ & $p \to q$ \\
\midrule
T & T & T \\
T & F & F \\
F & T & T \\
F & F & T \\
\bottomrule
\end{tabular}
\end{center}

Now let's show that with just or, and, and not we can replicate the implication operator:

\begin{center}
\begin{tabular}{cc|c}
\toprule
$p$ & $q$ & $(\lnot p) \vee q$ \\
\midrule
T & T & T \\
T & F & F \\
F & T & T \\
F & F & T \\
\bottomrule
\end{tabular}
\end{center}

Solving these little puzzles is the point of the homework this week.

\section{From truth table back to formula: a worked example}

Everything we've done so far goes in one direction: given a formula, compute its truth table. But in practice---especially when you're designing a digital circuit or implementing a controller that has to behave a certain way on every possible input---you often start with a truth table (a specification of what the output \emph{should} be) and need to work backward to a formula that realizes it.

Here's the worked example. Suppose you want a formula $\varphi(p, q, r)$ that is true exactly on the following rows:

\begin{center}
\begin{tabular}{ccc|c}
\toprule
$p$ & $q$ & $r$ & $\varphi(p,q,r)$ \\
\midrule
T & T & T & T \\
T & T & F & F \\
T & F & T & T \\
T & F & F & F \\
F & T & T & F \\
F & T & F & F \\
F & F & T & T \\
F & F & F & F \\
\bottomrule
\end{tabular}
\end{center}

\emph{Step 1.} Find every row where the output is T. Here there are three: $(T,T,T)$, $(T,F,T)$, and $(F,F,T)$.

\emph{Step 2.} For each such row, write a conjunction (``and'') that is true exactly on that row and false on every other row. The trick is to use the variable if the row has T for it, and its negation if the row has F:
\begin{itemize}
  \item Row $(T, T, T)$: $p \wedge q \wedge r$
  \item Row $(T, F, T)$: $p \wedge \lnot q \wedge r$
  \item Row $(F, F, T)$: $\lnot p \wedge \lnot q \wedge r$
\end{itemize}
Each of these little formulas is true on exactly one row of the truth table---check one of them if you like.

\emph{Step 3.} Take the \emph{disjunction} (``or'') of all those conjunctions:
\[
  \varphi(p, q, r) \;\equiv\; (p \wedge q \wedge r) \;\vee\; (p \wedge \lnot q \wedge r) \;\vee\; (\lnot p \wedge \lnot q \wedge r).
\]
This formula is true exactly when at least one of the three conjunctions is true, which is exactly on the three rows we wanted.

The form we just built---``or of ands, with each and being a combination of variables and their negations''---has a name: \emph{disjunctive normal form}\index{disjunctive normal form (DNF)}, or DNF. Every truth table can be converted to a DNF formula by this recipe. So propositional logic can express any truth function: given any desired behavior as a truth table, the DNF recipe gives you a formula that realizes it. That's actually a small theorem---\emph{every} boolean function on $n$ inputs is expressible in propositional logic---and it is half of why $\{\wedge, \vee, \lnot\}$ is functionally complete. (The other half is that any formula you can write in propositional logic \emph{is} such a boolean function, which is just what the truth-table semantics says.)

A side note: if you stare at the DNF formula for $\varphi$ above, you might notice that $r$ appears in every conjunction. That's no accident---looking at the original truth table, every T row has $r = T$. So we could have saved work by just writing $\varphi \equiv r \wedge \psi(p, q)$ for some simpler $\psi$. Simplifying a DNF formula is what a lot of digital-logic coursework is about, and it's the reason engineers care about tools like Karnaugh maps.

\section{Common mistakes}

Here are the mistakes students make most often on this module's material. Watch for them on the assignment.

\begin{enumerate}
  \item \textbf{Confusing implication with biconditional.} $p \to q$ is \emph{not} the same as $p \leftrightarrow q$. In particular, $p \to q$ is true when $p$ is false, regardless of what $q$ is; $p \leftrightarrow q$ is only true when $p$ and $q$ agree. If the English reads ``if and only if,'' you want $\leftrightarrow$; if the English reads ``if $\ldots$ then'' or ``$\ldots$ implies $\ldots$,'' you want $\to$.
  \item \textbf{Vacuous truth feels wrong.} The rows where $p$ is false make $p \to q$ come out true, and students often want those rows to be ``false'' or ``undefined.'' They aren't. $F \to q$ is \emph{true} by definition, because from a false premise we never make a false promise. If you can't accept this at gut level yet, at least accept it at definition level, and move on.
  \item \textbf{Operator precedence bugs.} Without parentheses, $\lnot p \wedge q$ means $(\lnot p) \wedge q$, \emph{not} $\lnot(p \wedge q)$. And $p \vee q \wedge r$ means $p \vee (q \wedge r)$, not $(p \vee q) \wedge r$. When in doubt, parenthesize explicitly; there's no points off for being unambiguous.
  \item \textbf{Inclusive vs.\ exclusive or.} Logical $\vee$ is \emph{inclusive}: $T \vee T = T$. Everyday English ``or'' is often exclusive---``coffee or tea'' usually means one or the other, not both. When you're translating English into propositional logic, be deliberate: did the sentence mean $\vee$ or $\oplus$?
  \item \textbf{Half-applied De Morgan.} Students often start pushing a negation inside and forget to flip the connective. $\lnot(p \wedge q)$ is $\lnot p \vee \lnot q$, \emph{not} $\lnot p \wedge \lnot q$. The connective flips: $\wedge \leftrightarrow \vee$.
\end{enumerate}

\section{Summary and looking ahead}

\begin{keytakeaway}
  \item \emph{Propositional logic} is defined by a BNF grammar over $T$, $F$, variables, $\wedge$, $\vee$, $\lnot$, $\to$. This is the first logic we define using the same ``syntax + semantics'' pattern we'll reuse for every logic that follows.
  \item The \emph{semantics} is given by \emph{truth tables}---a formula is a boolean-valued function of its variables.
  \item A \emph{tautology} is true under every assignment; a \emph{contradiction} is false under every assignment; everything else is \emph{contingent}. Two formulas are \emph{logically equivalent} iff they have the same truth table.
  \item De Morgan's laws push $\lnot$ across $\wedge$ and $\vee$, flipping the connective. This is the prototype for the quantifier negation rules in Module 4 ($\lnot\forall = \exists\lnot$, $\lnot\exists = \forall\lnot$).
  \item $\{\wedge, \vee, \lnot\}$ is \emph{functionally complete}---every boolean function can be written as a formula in these connectives. The DNF recipe (one conjunction per T row of the truth table) is the constructive proof.
\end{keytakeaway}

\begin{summarybox}
\textbf{What we built this module.} We have a full description of a logic: a grammar for writing formulas, a semantics for evaluating them, and enough machinery to check whether two formulas say the same thing. The truth-table method gives us a decision procedure---a mechanical way to answer any question about propositional-logic semantics---which is a rare luxury that won't survive the jump to predicate logic in Module 4. DNF is the bridge to circuit design in Module 3.

\medskip
\textbf{Looking ahead.} Module 3 asks a deeper question: can we \emph{prove} tautologies without computing truth tables? The answer is yes---we build a \emph{proof system} with introduction and elimination rules for each connective, and show those rules are sound against the truth-table semantics. The Gentzen intermezzo after Module 3 rewrites the same rules in the tree-shaped notation used by logicians and programming-language theorists, which will feel familiar from the DNF construction here.
\end{summarybox}

\section{Exercises}

These are organized into \textbf{warm-up} (mechanical practice on the definitions), \textbf{standard} (typical assignment-style problems), \textbf{challenge} (synthesis across the module), and \textbf{connection} (programming and intermezzo pointers). Solutions are in the appendix.

\subsection*{Warm-up}

\textbf{Exercise 1.}\label{ex:m2:1} Build the truth table for $(p \to q) \wedge (q \to p)$. This combination will come up later when we define the biconditional ($\leftrightarrow$)---you'll see why.

\bigskip
\textbf{Exercise 2.}\label{ex:m2:2} Prove the second De Morgan's law $\lnot(p \vee q) \equiv (\lnot p) \wedge (\lnot q)$ by computing both truth tables.

\bigskip
\textbf{Exercise 3.}\label{ex:m2:3} How many \emph{unary} boolean operations are there? List them all as truth tables, and for each one, give a short English name or description.

\bigskip
\textbf{Exercise 4.}\label{ex:m2:4} Classify each of the following formulas as a \textbf{tautology}, a \textbf{contradiction}, or \textbf{contingent} (neither always true nor always false). Do it by computing the truth table.
\begin{enumerate}[label=(\alph*)]
  \item $p \vee \lnot p$
  \item $p \wedge \lnot p$
  \item $p \to (q \to p)$
  \item $(p \to q) \to p$
  \item $(p \wedge q) \to (p \vee q)$
\end{enumerate}

\bigskip
\textbf{Exercise 5.}\label{ex:m2:5} For each of the following unparenthesized expressions, add parentheses to show how the precedence rules ($\lnot$ tighter than $\wedge$ tighter than $\vee$ tighter than $\to$) group them. Then evaluate the formula when $p = T$, $q = F$, $r = T$.
\begin{enumerate}[label=(\alph*)]
  \item $\lnot p \wedge q \vee r$
  \item $p \to q \to r$ \quad (Hint: $\to$ is conventionally \emph{right-associative}, so this groups as $p \to (q \to r)$, not $(p \to q) \to r$. Check what difference it makes.)
  \item $\lnot p \vee q \wedge \lnot r$
\end{enumerate}

\subsection*{Standard}

\textbf{Exercise 6.}\label{ex:m2:6} Show that XOR ($p \oplus q$) can be expressed using only $\wedge$, $\vee$, and $\lnot$. Hint: compare truth tables and use the DNF recipe from the worked example.

\bigskip
\textbf{Exercise 7.}\label{ex:m2:7} Prove the equivalence $p \to (q \wedge r) \;\equiv\; (p \to q) \wedge (p \to r)$ by truth table. (This is sometimes called ``implication distributes over conjunction on the right.'' It means that to prove a conjunction from a premise, it suffices to prove each conjunct separately.)

\bigskip
\textbf{Exercise 8.}\label{ex:m2:8} Show how to build each of $\lnot$, $\wedge$, and $\vee$ using only the NAND operator. That is:
\begin{enumerate}[label=(\alph*)]
  \item Find a formula using only NAND (and the variable $p$) whose truth table matches $\lnot p$.
  \item Find a formula using only NAND whose truth table matches $p \wedge q$. (Hint: NAND is ``not-and,'' so you can build AND by negating NAND, and you already know how to build NOT from NAND by part (a).)
  \item Find a formula using only NAND whose truth table matches $p \vee q$. (Hint: De Morgan. $p \vee q \equiv \lnot(\lnot p \wedge \lnot q)$, and you can build all three ingredients from NAND.)
\end{enumerate}
This shows that NAND alone is functionally complete, which is why circuit designers love it: you can build an entire CPU out of NAND gates.

\bigskip
\textbf{Exercise 9.}\label{ex:m2:9} A three-input boolean function $f(x, y, z)$ outputs T on exactly the input rows $(T, F, F)$, $(F, T, F)$, and $(F, F, T)$ (the ``exactly-one'' function on three inputs). Use the DNF recipe from the worked example to write a propositional formula for $f$. Your formula should have three disjuncts, each with three literals.

\subsection*{Challenge}

\textbf{Exercise 10.}\label{ex:m2:10} Show that $\{\wedge, \vee\}$---conjunction and disjunction, without negation---is \emph{not} functionally complete. Hint: call a function $f$ \emph{monotone} if flipping any input from $F$ to $T$ never flips an output from $T$ to $F$. (a) Show by induction on the structure of formulas that any formula built from just $\wedge$ and $\vee$ represents a monotone function. (b) Exhibit a boolean function that is not monotone. (c) Conclude.

\bigskip
\textbf{Exercise 11.}\label{ex:m2:11} Prove: a formula $\varphi$ is a tautology if and only if $\lnot \varphi$ is a contradiction. (Do this both ways. Assume $\varphi$ is a tautology and show $\lnot \varphi$ is a contradiction; then assume $\lnot \varphi$ is a contradiction and show $\varphi$ is a tautology.) This is the definitional glue that lets proof by contradiction work---more on that in Module 7.

\bigskip
\textbf{Exercise 12.}\label{ex:m2:12} Of the 16 binary boolean operations:
\begin{enumerate}[label=(\alph*)]
  \item How many depend on \emph{neither} of their inputs (i.e., the output is constant regardless of the inputs)? List them.
  \item How many depend on \emph{only one} of their inputs (so they ignore the other)? List them.
  \item How many depend on \emph{both} inputs?
\end{enumerate}
(Your three counts should add up to 16.)

\subsection*{Connection}

\textbf{Exercise 13.}\label{ex:m2:13} (Programming.) Write a Python function \texttt{is\_tautology(f, n)} that takes a function \texttt{f} of $n$ boolean arguments and returns \texttt{True} if \texttt{f} is a tautology (i.e., returns \texttt{True} on every input) and \texttt{False} otherwise. Test it on \texttt{lambda p, q: (not p) or q} (which is logically equivalent to $p \to q$, so \emph{not} a tautology, only a contingent formula) and on \texttt{lambda p, q: p or (not p)} (which \emph{is} a tautology). Hint: use \texttt{itertools.product([False, True], repeat=n)} to enumerate all assignments.

\bigskip
\textbf{Exercise 14.}\label{ex:m2:14} (Intermezzo pointer.) The discussion in this module of why the law of excluded middle ($p \vee \lnot p$) is controversial points forward to the ``Proofs Are Programs'' intermezzo (Curry--Howard), which you can read after Module 7. As a preview: try to write a \emph{program} that, given a predicate $P$ on natural numbers, produces either a natural number that satisfies $P$ or a proof that no such number exists. What would the input and output types of that program have to be, and why is it not obvious how to implement it? (You don't need to write actual code---just think about the type signature.)

\end{document}
